\documentclass{article}
\usepackage[utf8]{inputenc}
\usepackage{geometry}
\usepackage{titling}
\usepackage{titlesec}
\usepackage{enumitem}
\usepackage{hyperref}

\geometry{
 a4paper,
 total={170mm,257mm},
 left=20mm,
 top=20mm,
}

\title{\textbf{Research Framework: AI Voice Detection}}
\author{Research Planning Document}
\date{\today}

\begin{document}

\maketitle

\section{Research Objectives}

The primary objective is to develop a real-time detection system for AI-generated voices, with particular emphasis on deepfake detection. The research must distinguish between detection (identifying synthetic speech) and authentication (verifying speaker identity), as both require different approaches.

\section{Core Research Areas}

Two foundational areas require thorough investigation before addressing the hardware evaluation question:

\subsection{Biology and Physics of Speech}

Understanding natural speech production mechanisms is essential for identifying what neural TTS systems cannot or do not capture. Key aspects include:

\begin{itemize}
    \item \textbf{Source-Filter Theory}: The glottal source (vocal cord vibration) and vocal tract filter (resonances/formants) and their coupling dynamics
    \item \textbf{Vocal Cord Physics}: Biomechanics of vocal fold vibration, subglottal pressure dynamics, and physical constraints
    \item \textbf{Frequency Range Analysis}: Sub-50Hz (vocal cord micro-movements), 50-500Hz (fundamental frequency), 500-8000Hz (formants), and $>$20kHz (ultrasonic content)
    \item \textbf{Temporal Dynamics}: Micro-timing variations ($<$1ms), jitter, shimmer, and natural temporal irregularities
\end{itemize}

\subsection{TTS System Internals}

Understanding neural TTS architectures from first principles rather than analyzing their outputs alone:

\begin{itemize}
    \item \textbf{Architecture Analysis}: Autoregressive vs. parallel generation, encoder-decoder structures, attention mechanisms, receptive field limitations
    \item \textbf{Training Objectives}: Likelihood-based training, adversarial training, reconstruction losses, and their systematic biases
    \item \textbf{Latent Space Structure}: Variational autoencoder latent spaces, flow-based transformations, manifold structure
    \item \textbf{Information-Theoretic Properties}: Entropy characteristics, mutual information, statistical regularities
    \item \textbf{Optimization Artifacts}: Gradient descent path effects, convergence behavior, regularization impacts
\end{itemize}

\section{Detection Methodology}

Current detection approaches rely on artifact identification (phase patterns, spectral smoothness), which are easily patchable as TTS systems evolve. Alternative approaches include:

\subsection{Algorithmic Vulnerability Analysis}

Rather than detecting surface artifacts, analyze TTS systems as algorithms to identify fundamental vulnerabilities:
\begin{itemize}
    \item Neural network architecture constraints
    \item Training objective function biases
    \item Latent space regularity
    \item Model capacity limitations
\end{itemize}

\subsection{Universal Speech Laws}

Detect violations of acoustic physics that all neural TTS systems must obey:
\begin{itemize}
    \item Source-filter theory constraints
    \item Vocal tract physics limitations
    \item Prosodic universals
    \item Temporal dependencies
\end{itemize}

\subsection{Extended Frequency Analysis}

Analyze frequency ranges that standard TTS systems (typically trained on 16-24kHz) may not capture:
\begin{itemize}
    \item Sub-50Hz components
    \item Ultrasonic content ($>$20kHz)
    \item Micro-dynamics ($<$1ms)
    \item Sub-harmonic components
\end{itemize}

\section{Hardware Evaluation Question}

The question of whether microphone quality is a detection bottleneck can only be answered after establishing:
\begin{enumerate}
    \item What physical characteristics define natural speech
    \item What fundamental limitations TTS architectures have
    \item What features require better hardware to capture
\end{enumerate}

This foundational understanding is necessary before evaluating hardware impact.

\section{Authentication vs. Detection}

\subsection{Detection}

Primary focus: Identifying whether audio is AI-generated or human-produced in real-time.

\subsection{Authentication}

Secondary consideration: Verifying speaker identity, which may involve:
\begin{itemize}
    \item Voice signature capture (spectral envelope, prosody, timing, F0 contour)
    \item Long-term patterns (speaking rate, pause patterns, phoneme transitions)
    \item Micro-features (jitter, shimmer, breath patterns)
    \item Liveness detection (challenge-response, semantic inconsistency, emotional response)
\end{itemize}

Authentication is distinct from detection and requires different approaches.

\section{Research Priorities}

\begin{enumerate}
    \item Understand biology and physics of speech production
    \item Understand TTS system internals and architecture
    \item Develop detection methods based on fundamental understanding
    \item Evaluate hardware impact with evidence-based conclusions
\end{enumerate}

\end{document}
